% Developer Tooling & Agile section

% 1 – Git Core Concepts
\QuestionSlide[\CategoryBadge[ToolingColor!20]{Tooling \& Agile}]{* What are the core Git concepts every developer must know?}

\begin{frame}[fragile]
  \frametitle{\AnswerTitle[\CategoryBadge[ToolingColor!20]{Tooling \& Agile}]{Core Git concepts?}}

  {\footnotesize
  \begin{itemize}
    \item \textbf{Repository} – the \texttt{.git} directory storing history as a DAG of commits.
    \item \textbf{Working tree / Index (staging area)} – prepare commits selectively.
    \item \textbf{Branch} – lightweight movable pointer to a commit.
    \item \textbf{Merge vs Rebase} – merge preserves history; rebase rewrites it to a linear sequence.
    \item \textbf{Cherry-pick} – apply a specific commit to another branch.
    \item \textbf{Stash} – save work-in-progress without committing.
  \end{itemize}
  }

  \begin{minted}[fontsize=\scriptsize]{bash}
git switch -c feature/auth    # new branch
git add src/Auth.cs           # stage specific file
git commit -m "feat: add JWT auth"
git rebase main               # linear history, clean PRs
git stash && git switch main  # stash WIP, switch context
  \end{minted}
\end{frame}

% 2 – Git Branching Strategies
\QuestionSlide[\CategoryBadge[ToolingColor!20]{Tooling \& Agile}]{** What are GitFlow and Trunk-Based Development?}

\begin{frame}[fragile]
  \frametitle{\AnswerTitle[\CategoryBadge[ToolingColor!20]{Tooling \& Agile}]{GitFlow vs Trunk-Based Development?}}

  {\footnotesize
  \begin{tabular}{ll}
    \textbf{GitFlow} & \textbf{Trunk-Based Development} \\
    Long-lived \texttt{develop} + \texttt{main} & Single \texttt{main} (trunk) \\
    \texttt{feature/}, \texttt{release/}, \texttt{hotfix/} & Short-lived feature branches ($<$1 day) \\
    Slower release cadence & Continuous delivery, multiple deploys/day \\
    Complex merge overhead & Feature flags hide in-progress work \\
    Suited for scheduled releases & Suited for CI/CD and microservices \\
  \end{tabular}
  }
\end{frame}

% 3 – Azure DevOps Boards
\QuestionSlide[\CategoryBadge[ToolingColor!20]{Tooling \& Agile}]{* What is Azure DevOps Boards?}

\begin{frame}[fragile]
  \frametitle{\AnswerTitle[\CategoryBadge[ToolingColor!20]{Tooling \& Agile}]{What is Azure DevOps Boards?}}

  {\footnotesize
  \textbf{Azure DevOps Boards} is Microsoft's Agile planning tool within Azure DevOps. Key work item types: \textbf{Epic} $\to$ \textbf{Feature} $\to$ \textbf{User Story} $\to$ \textbf{Task} / \textbf{Bug}. It offers Kanban boards, Sprint backlogs, burndown charts, and \textbf{Delivery Plans} for cross-team scheduling. Integrates with GitHub PRs and Azure Pipelines for end-to-end traceability.
  }
\end{frame}

% 4 – Jira
\QuestionSlide[\CategoryBadge[ToolingColor!20]{Tooling \& Agile}]{* What is Jira and when would you choose it over Azure DevOps Boards?}

\begin{frame}[fragile]
  \frametitle{\AnswerTitle[\CategoryBadge[ToolingColor!20]{Tooling \& Agile}]{Jira vs Azure DevOps Boards?}}

  {\footnotesize
  \textbf{Jira} (Atlassian) is the most widely adopted issue tracker and Agile board tool. Strengths: highly customisable workflows, 3000+ integrations (Confluence, Bitbucket, GitHub), JQL (query language), and Jira Product Discovery for roadmaps. Choose Jira if you use the broader Atlassian stack or need complex cross-project workflows. Choose Azure DevOps Boards for deep Microsoft/Azure pipeline integration.
  }
\end{frame}

% 5 – Agile / Scrum
\QuestionSlide[\CategoryBadge[ToolingColor!20]{Tooling \& Agile}]{* What is Scrum and its key ceremonies?}

\begin{frame}[fragile]
  \frametitle{\AnswerTitle[\CategoryBadge[ToolingColor!20]{Tooling \& Agile}]{What is Scrum?}}

  {\footnotesize
  \textbf{Scrum} is an Agile framework where work is delivered in fixed \textbf{Sprints} (1–4 weeks). Roles: \textbf{Product Owner} (owns backlog), \textbf{Scrum Master} (removes blockers), \textbf{Development Team}. Key ceremonies:
  \begin{itemize}
    \item \textbf{Sprint Planning} – select and break down stories.
    \item \textbf{Daily Standup} – 15-min sync on blockers.
    \item \textbf{Sprint Review} – demo to stakeholders.
    \item \textbf{Retrospective} – process improvement.
  \end{itemize}
  }
\end{frame}

% 6 – Power BI
\QuestionSlide[\CategoryBadge[ToolingColor!20]{Tooling \& Agile}]{* What is Power BI?}

\begin{frame}[fragile]
  \frametitle{\AnswerTitle[\CategoryBadge[ToolingColor!20]{Tooling \& Agile}]{What is Power BI?}}

  {\footnotesize
  \textbf{Power BI} is Microsoft's business intelligence platform for building interactive dashboards and reports. It ingests data from SQL Server, Azure Synapse, Dataverse, Excel, REST APIs, and many connectors. \textbf{Power BI Desktop} (Windows) authors reports; \textbf{Power BI Service} (cloud) shares and schedules refreshes. \textbf{DAX} is the formula language for calculated measures and KPIs.
  }
\end{frame}

% 7 – Definition of Done
\QuestionSlide[\CategoryBadge[ToolingColor!20]{Tooling \& Agile}]{** What is the Definition of Done (DoD) in Scrum?}

\begin{frame}[fragile]
  \frametitle{\AnswerTitle[\CategoryBadge[ToolingColor!20]{Tooling \& Agile}]{What is the Definition of Done?}}

  {\footnotesize
  The \textbf{Definition of Done (DoD)} is the team's shared checklist that a Story must pass before it is considered complete. A typical DoD for a .NET team:
  \begin{itemize}
    \item Code reviewed and approved.
    \item Unit and integration tests written and green.
    \item CI pipeline passes (build, test, static analysis).
    \item Feature deployed to staging and smoke-tested.
    \item Acceptance criteria verified by Product Owner.
    \item Documentation updated (API spec, ADR).
  \end{itemize}
  A weak DoD leads to hidden technical debt and unstable releases.
  }
\end{frame}

% 8 – LaTeX
\QuestionSlide[\CategoryBadge[ToolingColor!20]{Tooling \& Agile}]{* What is LaTeX and when should a developer use it?}

\begin{frame}[fragile]
  \frametitle{\AnswerTitle[\CategoryBadge[ToolingColor!20]{Tooling \& Agile}]{What is LaTeX?}}

  {\footnotesize
  \textbf{LaTeX} is a document typesetting system where content and layout are expressed as markup (\texttt{.tex} files). It produces professional-quality PDFs with precise control over typography, mathematical formulas, code listings, and references. Ideal for: academic papers, technical reports, CV/slides (\texttt{beamer}), and this very flashcard deck. Store \texttt{.tex} files in Git for version-controlled documents.
  }

  \begin{minted}[fontsize=\scriptsize]{bash}
# Compile to PDF
latexmk -xelatex -shell-escape main.tex

# Typical project structure
# main.tex     – entry point, \input{} sections
# sections/    – chapter/section .tex files
# figures/     – vector graphics (PDF, TikZ)
  \end{minted}
\end{frame}
